\section{P04 - Barème - Types construits}

\subsection{Objectifs spécifiques}

\begin{itemize}
\tightlist
\item
  Objectif O1 - Identifier précisément la représentation ou la structure
  en jeu.
\item
  Objectif O2 - Appliquer une méthode disciplinaire complète.
\item
  Objectif O3 - Justifier le résultat sur un cas différent.
\item
  Objectif O4 - Contrôler un cas limite et corriger une erreur observée.
\end{itemize}

\subsection{Capacités officielles atomiques}

\begin{itemize}
\tightlist
\item
  P-DATA-CONSTR-02A
\end{itemize}

\subsection{Prérequis}

\begin{itemize}
\tightlist
\item
  Reconnaître une consigne liée à tuple.
\item
  Distinguer donnée, méthode et conclusion dans le thème Tuples, listes,
  dictionnaires.
\item
  Rédiger une justification courte en utilisant le vocabulaire du
  programme.
\item
  Contrôler une réponse par un cas limite ou un contre-exemple
  explicite.
\end{itemize}

\subsection{Séance(s) correspondante(s)}

\begin{itemize}
\tightlist
\item
  P04-S1 à P04-S7 : support rattaché aux séances prêtes de la
  progression.
\end{itemize}

\subsection{Situation-problème concrète}

Une station météo stocke des coordonnées fixes, des relevés horaires
modifiables et des mesures accessibles par nom.

\subsection{Activité d’entrée}

\begin{enumerate}
\def\labelenumi{\arabic{enumi}.}
\tightlist
\item
  Identifier ce qui doit rester immuable dans un tuple.
\item
  Modifier une liste de températures.
\item
  Lire une clé dans un dictionnaire de station.
\item
  Décrire ce qui se passe avec une liste vide.
\end{enumerate}

\subsubsection{Barème question 1}

\begin{itemize}
\tightlist
\item
  1 point : identifier correctement tuple de coordonnées et la donnée
  \texttt{(36.8,\ 10.2)}.
\item
  1 point : appliquer la méthode « lire sans modifier et nommer latitude
  puis longitude ».
\item
  1 point : obtenir coordonnées conservées.
\item
  1 point : contrôler le cas limite « tentative de modification
  interdite » ou expliquer pourquoi il ne s'applique pas.
\item
  Retrait possible : confusion avec l'erreur fréquente EF1 - Modifier un
  tuple comme une liste. \#\#\# Barème question 2
\item
  1 point : identifier correctement liste de relevés et la donnée
  \texttt{{[}18,\ 20,\ 19{]}}.
\item
  1 point : appliquer la méthode « parcourir les valeurs et calculer une
  moyenne ».
\item
  1 point : obtenir \texttt{19}.
\item
  1 point : contrôler le cas limite « liste vide » ou expliquer pourquoi
  il ne s'applique pas.
\item
  Retrait possible : confusion avec l'erreur fréquente EF2 - Parcourir
  les indices quand les valeurs suffisent. \#\#\# Barème question 3
\item
  1 point : identifier correctement dictionnaire et la donnée
  \texttt{\{"temperature":\ 21,\ "vent":\ 12\}}.
\item
  1 point : appliquer la méthode « tester la présence de la clé avant
  lecture ».
\item
  1 point : obtenir \texttt{21} pour \texttt{temperature}.
\item
  1 point : contrôler le cas limite « clé absente » ou expliquer
  pourquoi il ne s'applique pas.
\item
  Retrait possible : confusion avec l'erreur fréquente EF3 - Accéder à
  une clé sans vérifier sa présence. \#\#\# Barème question 4
\item
  1 point : identifier correctement copie de liste et la donnée
  \texttt{{[}{[}1{]},\ {[}2{]}{]}}.
\item
  1 point : appliquer la méthode « distinguer copie superficielle et
  copie indépendante ».
\item
  1 point : obtenir modification locale contrôlée.
\item
  1 point : contrôler le cas limite « liste imbriquée » ou expliquer
  pourquoi il ne s'applique pas.
\item
  Retrait possible : confusion avec l'erreur fréquente EF4 - Copier une
  liste imbriquée seulement au premier niveau. \#\# Exercices numérotés
\item
  Exercice 1 : résoudre tuple de coordonnées avec \texttt{(36.8,\ 10.2)}
  ; attendu : coordonnées conservées.
\item
  Exercice 2 : expliquer liste de relevés à partir de
  \texttt{{[}18,\ 20,\ 19{]}} ; attendu : \texttt{19}.
\item
  Exercice 3 : comparer dictionnaire avec
  \texttt{\{"temperature":\ 21,\ "vent":\ 12\}} ; attendu : \texttt{21}
  pour \texttt{temperature}.
\item
  Exercice 4 : corriger copie de liste pour
  \texttt{{[}{[}1{]},\ {[}2{]}{]}} ; attendu : modification locale
  contrôlée.
\item
  Exercice 5 : tester un cas limite lié à tentative de modification
  interdite ; attendu : le comportement de tuple de coordonnées est
  contrôlé.
\item
  Exercice 6 : classer deux méthodes possibles pour liste de relevés ;
  attendu : la méthode robuste est choisie et justifiée.
\item
  Exercice 7 : justifier un transfert qui utilise dictionnaire avec une
  donnée nouvelle ; attendu : la justification reste valable sur le
  nouveau cas.
\item
  Exercice 8 : étendre un énoncé volontairement erroné sur copie de
  liste ; attendu : l'erreur est localisée puis réparée.
\end{itemize}

\subsection{Corrigés complets des exercices du cours}

\begin{itemize}
\tightlist
\item
  Corrigé exercice 1 : méthode : identifier \texttt{(36.8,\ 10.2)},
  appliquer la méthode « lire sans modifier et nommer latitude puis
  longitude », puis écrire coordonnées conservées ; résultat :
  coordonnées conservées ; contrôle : faire apparaître le contrôle «
  tentative de modification interdite ».
\item
  Corrigé exercice 2 : méthode : expliciter chaque étape de parcourir
  les valeurs et calculer une moyenne avant de conclure par \texttt{19}
  ; résultat : \texttt{19} ; contrôle : rédiger la méthode avant le
  résultat.
\item
  Corrigé exercice 3 : méthode : comparer la donnée avec le cas limite «
  clé absente » et valider \texttt{21} pour \texttt{temperature} ;
  résultat : \texttt{21} pour \texttt{temperature} ; contrôle : comparer
  avec le cas « clé absente ».
\item
  Corrigé exercice 4 : méthode : isoler l'erreur fréquente « Copier une
  liste imbriquée seulement au premier niveau. » puis reprendre la
  procédure correcte ; résultat : modification locale contrôlée ;
  contrôle : corriger l'erreur « Copier une liste imbriquée seulement au
  premier niveau. ».
\item
  Corrigé exercice 5 : méthode : identifier \texttt{(36.8,\ 10.2)},
  appliquer la méthode « lire sans modifier et nommer latitude puis
  longitude », puis écrire coordonnées conservées ; résultat : le
  comportement de tuple de coordonnées est contrôlé ; contrôle : nommer
  la donnée minimale et la conclusion.
\item
  Corrigé exercice 6 : méthode : expliciter chaque étape de parcourir
  les valeurs et calculer une moyenne avant de conclure par \texttt{19}
  ; résultat : la méthode robuste est choisie et justifiée ; contrôle :
  identifier pourquoi « Parcourir les indices quand les valeurs
  suffisent. » est une erreur.
\item
  Corrigé exercice 7 : méthode : comparer la donnée avec le cas limite «
  clé absente » et valider \texttt{21} pour \texttt{temperature} ;
  résultat : la justification reste valable sur le nouveau cas ;
  contrôle : inclure une étape calculable par un pair.
\item
  Corrigé exercice 8 : méthode : isoler l'erreur fréquente « Copier une
  liste imbriquée seulement au premier niveau. » puis reprendre la
  procédure correcte ; résultat : l'erreur est localisée puis réparée ;
  contrôle : proposer une activité corrective inspirée de « Modifier une
  sous-liste et observer l'effet sur la copie. ».
\end{itemize}

\subsection{Erreurs fréquentes}

\begin{itemize}
\tightlist
\item
  Erreur fréquente EF1 - Modifier un tuple comme une liste.
\item
  Erreur fréquente EF2 - Parcourir les indices quand les valeurs
  suffisent.
\item
  Erreur fréquente EF3 - Accéder à une clé sans vérifier sa présence.
\item
  Erreur fréquente EF4 - Copier une liste imbriquée seulement au premier
  niveau.
\end{itemize}

\subsection{Remédiation ciblée}

\begin{itemize}
\tightlist
\item
  Activité corrective EF1 : Identifier mutabilité et usage avant
  d'écrire une affectation.
\item
  Activité corrective EF2 : Écrire deux boucles, avec indices puis avec
  valeurs, et comparer.
\item
  Activité corrective EF3 : Tester \texttt{cle\ in\ dictionnaire} avant
  la lecture.
\item
  Activité corrective EF4 : Modifier une sous-liste et observer l'effet
  sur la copie.
\end{itemize}

\subsection{Différenciation}

\begin{itemize}
\tightlist
\item
  Socle : traiter \texttt{(36.8,\ 10.2)} avec une fiche méthode fournie.
\item
  Standard : traiter \texttt{{[}18,\ 20,\ 19{]}} en rédigeant la
  justification complète.
\item
  Expert : inventer un cas limite lié à « clé absente » et expliquer le
  comportement attendu.
\end{itemize}

\subsection{Critères de réussite}

\begin{itemize}
\tightlist
\item
  La capacité officielle est citée dans la copie.
\item
  La méthode contient au moins une étape vérifiable par un pair.
\item
  Le cas limite est discuté avec une donnée concrète.
\item
  La correction explique quelle erreur fréquente est évitée.
\end{itemize}

\subsubsection{Barème question 5}

\begin{itemize}
\tightlist
\item
  1 point : identifier le tuple comme coordonnée immuable.
\item
  1 point : identifier le dictionnaire comme accès par clé.
\item
  1 point : identifier la liste comme collection parcourable.
\item
  1 point : calculer le milieu \texttt{(3.0,\ 7.0)} avec formule
  correcte.
\end{itemize}
